\chapter{The \texttt{.deppy} Sidecar Format}
\label{ch:deppy}

A sidecar file (\texttt{*.deppy}) is the user-facing interface to
the \Deppy{} pipeline.  It declaratively records, for one Python
library:

\begin{itemize}[leftmargin=1.7em]
  \item a small number of \emph{foundations} (mathematical facts the
    user is willing to admit);
  \item one or more \emph{specs} (interface contracts on classes /
    functions);
  \item \emph{axioms} attached to specs (claims about library
    behaviour);
  \item \emph{lemmas} (intermediate facts the user wants to prove
    once and re-use);
  \item \emph{verifies} (the unit of verification: function +
    property + proof script);
  \item \emph{proofs} (legacy form: name + by-clause selecting a
    derivation strategy);
  \item top-level \emph{code\_types} (signatures for helper
    functions PSDL references);
  \item top-level \emph{lean\_imports} / \emph{lean\_preamble} (raw
    \Lean{} text injected verbatim).
\end{itemize}

The file is read by \texttt{SidecarModule.from\_file} (file:
\texttt{deppy/proofs/sidecar.py}) and parsed into a
\texttt{SidecarModule} containing
\texttt{FoundationDecl} / \texttt{SpecDecl} / \texttt{AxiomDecl} /
\texttt{LemmaDecl} / \texttt{VerifyDecl} / \texttt{ProofDecl}
dataclasses.  The pipeline then operates exclusively on this
\texttt{SidecarModule}.

\section{File-level header}

Every sidecar starts with a module identifier:

\begin{lstlisting}[language=deppyfile,caption={Header of \texttt{sympy\_geometry.deppy}.}]
module: sympy
version: ">=1.12"
\end{lstlisting}

The \texttt{module:} field is the dotted name of the target Python
package, used by the runner (\texttt{deppy.lean.sidecar\_runner}) to
look up provenance information (version, source path) via
\texttt{importlib.import\_module}.  The \texttt{version:} field is
informational; the pipeline does not currently enforce it.  A future
extension could compare against \verb|module.__version__|.

\section{Foundations}
\label{sec:foundations}

A foundation is a generic mathematical fact admitted as part of the
trust surface.  Syntax:

\begin{lstlisting}[language=deppyfile]
foundation Real_sqrt_nonneg: "sqrt(x) >= 0 when x >= 0"
foundation Real_add_comm:    "a + b == b + a"
foundation Real_sub_sq_swap: "(a - b)**2 == (b - a)**2"
\end{lstlisting}

Each foundation has a name (a \Lean{}-safe identifier when
underscored) and a Python-syntax statement.  A trailing
\texttt{when ...} guard introduces a precondition.

The pipeline tries each foundation through \Zthree{}
(\Cref{ch:z3}); foundations involving \texttt{sqrt}, \texttt{sin},
\texttt{cos}, \dots fall through and stay admitted.  In the example
sympy file, \texttt{Real\_add\_comm} is provable by \texttt{omega}
in \Lean, so the certificate emits a real
\texttt{theorem Foundation\_Real\_add\_comm} with body
\texttt{by intros; omega}; \texttt{Real\_sqrt\_nonneg} stays
\texttt{axiom}-admitted because Z3 over \texttt{Int} cannot model
\texttt{sqrt}.

A foundation can also carry an explicit user-supplied \Lean{} proof:
\begin{lstlisting}[language=deppyfile]
foundation Real_my_fact:
  statement: "<python statement>"
  citation:  "Bourbaki, Algebra II, §3.4"
  lean_signature: "<optional Lean type override>"
  lean_proof: |
    intros
    omega
\end{lstlisting}
When \texttt{lean\_proof} is non-empty, the certificate emits a real
\texttt{theorem ... := by <lean\_proof>} instead of an
\texttt{axiom}; \Lean{} must elaborate the proof.

\section{Specs and axioms}
\label{sec:specs-axioms}

A \texttt{spec} attaches a contract bundle to a class or a function.
Inside it, \texttt{axiom} declarations attach \emph{claims about
behaviour} to that target.

\begin{lstlisting}[language=deppyfile,caption={Spec for
\texttt{sympy.geometry.point.Point}.}]
spec sympy.geometry.point.Point:
  guarantee: "Euclidean distance is a metric on R²; collinearity is the vanishing cross product"
  pure: true
  axiom dist_nonneg: "Point.distance(p, q) >= 0"
  axiom dist_zero_iff_equal: "(Point.distance(p, q) == 0) == (p == q)"
  axiom dist_symmetric: "Point.distance(p, q) == Point.distance(q, p)"
  axiom dist_triangle: "Point.distance(p, r) <= Point.distance(p, q) + Point.distance(q, r)"
  axiom dist_pythagoras: "Point.distance(p, q)**2 == (p.x - q.x)**2 + (p.y - q.y)**2"
  axiom collinear_iff_zero_cross: "is_collinear(A, B, C) == ((B.x-A.x)*(C.y-A.y) == (C.x-A.x)*(B.y-A.y))"
\end{lstlisting}

The fields:
\begin{description}
  \item[\texttt{guarantee:}] a free-text human-readable summary of the
  spec's contract.  Used in the certificate as a header comment.

  \item[\texttt{pure: true / false}] purity claim.  Pure functions
  inhabit the $\Box$-modality (no exceptions, no I/O, no mutation).

  \item[\texttt{requires:} / \texttt{ensures:}] (optional) Hoare-style
  pre/postconditions, in Python expression form.

  \item[\texttt{reads:} / \texttt{mutates:}] (optional) frame
  conditions for impure functions.

  \item[\texttt{axiom NAME: "...".}] a sub-declaration that becomes
  an \texttt{AxiomDecl} attached to the enclosing spec.
\end{description}

The axiom statement is a Python expression.  The spec's target
provides the namespace: an \texttt{axiom dist\_nonneg} on
\texttt{spec ...Point} compiles to a \Lean{} \texttt{Prop}
quantifying over the binders implied by the statement.

\section{Verify blocks}
\label{sec:verify-blocks}

A \texttt{verify} block is the unit of verification: it names a
function, a property, and a proof strategy.  It is the most direct
way to drive the kernel pipeline.  The full schema (file:
\texttt{deppy/proofs/sidecar.py}, dataclass \texttt{VerifyDecl},
lines 388--446):

\begin{lstlisting}[language=deppyfile,caption={Schema of a
\texttt{verify} block (all fields shown; only \texttt{function} and
\texttt{property} are mandatory).}]
verify <NAME>:
  function:  <dotted path to method>
  property:  "<Python predicate>"
  via:       foundation <foundation-name>     # OR axiom <axiom-name>
  fibration: isinstance <Type>                # optional
  cubical:   true | false                     # optional
  binders:   p: Point, q: Point               # optional, drives concrete Props
  requires:  "<precondition>"                 # optional
  ensures:   "<postcondition>"                # optional
  effect:    "pure"                           # optional
  decreases: "<measure expression>"           # optional
  z3_binders: x: int, y: int                  # optional
  cubical_dim: 0                              # optional
  cohomology_level: 0                         # optional
  expects_sha256: "<sha-256 hex>"             # optional, hash drift check
  expecting_counterexample: false             # optional
  by_lean: |                                  # optional, raw Lean tactic
    intros
    apply Foundation_X
    omega
  psdl_proof: |                               # optional, PSDL block
    if isinstance(other, Point):
        ...
\end{lstlisting}

\subsection{Mandatory fields}

\paragraph{\texttt{function}.}  A dotted path that
\texttt{importlib} can resolve down to a callable (or a
\texttt{property} descriptor whose \texttt{fget} is a callable).
The pipeline reads its source via \texttt{inspect.getsource} and
hashes it for the \texttt{source\_hash} field.

\paragraph{\texttt{property}.}  A Python expression in the function's
binder context.  The expression is parsed via
\texttt{ast.parse(\dots, mode="eval")}; identifiers that look like
free variables are quantified with the binder types from
\texttt{binders:}; identifiers that look like classes resolve via
the \texttt{code\_types:} table; identifiers that match the function's
own class invoke the method machinery from
\texttt{deppy.lean.sidecar\_mechanization}.

\subsection{Discharge strategy}

The \texttt{via:} clause selects a discharge strategy.  Currently
recognised forms:

\begin{itemize}[leftmargin=1.7em]
  \item \texttt{via: foundation NAME} --- cite the named foundation
    as the proof.
  \item \texttt{via: axiom NAME} --- cite a sidecar axiom.
  \item \texttt{via: chain F1, F2, F3} --- transitively compose the
    named axioms/foundations via \texttt{Eq.trans}.
  \item \texttt{via: rewrite F then cite A} --- rewrite under
    foundation $F$, then cite axiom $A$.
  \item \texttt{via: composition F1 \& F2 \& F3} --- and-intro a
    conjunction of named claims.
\end{itemize}

When neither \texttt{psdl\_proof} nor \texttt{by\_lean} is given, the
\texttt{via:} clause is the proof strategy and the kernel \emph{builds}
the appropriate \texttt{ProofTerm} (an \texttt{AxiomInvocation},
\texttt{Trans} chain, \texttt{Rewrite}, or \texttt{Cut}).

\subsection{Proof script fields}

The two proof script fields are mutually compatible:

\begin{description}
  \item[\texttt{psdl\_proof: |}] A multi-line Python-shaped tactic
  script.  Each statement compiles to one or more
  \texttt{ProofTerm}s and one or more \Lean{} tactic lines; the
  resulting \Lean{} block is wrapped in \texttt{by ...} for the
  generated theorem.  See \Cref{ch:psdl}.

  \item[\texttt{by\_lean: |}] Raw \Lean{} tactic text, inserted
  verbatim.  The user is responsible for the tactic body
  elaborating against the concrete \texttt{Prop}; the \Deppy{}
  pipeline only checks for SHA hash drift on the source.
\end{description}

When both are provided, \texttt{by\_lean} takes precedence in the
\Lean{} certificate but \texttt{psdl\_proof} still produces a
\texttt{ProofTerm} the kernel checks.  This is the primary way to
provide ``a proof in two notations'' for the same property.

\subsection{Optional fields driving concretisation}

The annotation-driven preamble (\Cref{ch:lean}) reads:

\begin{itemize}[leftmargin=1.7em]
  \item \texttt{binders: p: Point, q: Point} --- declares
    \texttt{axiom Point : Type}, then uses these as
    \texttt{intros p q} and as the \texttt{Verified\_X\_property}'s
    quantifiers.
  \item \texttt{requires: "p != q"} --- an extra hypothesis that the
    generated \texttt{Verified\_X\_property} body assumes.
  \item \texttt{ensures: "<expr>"} --- an extra clause appended to
    the property body.
  \item \texttt{effect: "pure"} or \texttt{"may\_raise TypeError"}
    --- recorded for the certificate as a comment; informs the
    fibration-descent pass.
  \item \texttt{decreases: "<measure>"} --- a well-founded measure
    used by the safety pipeline to discharge recursion-depth
    obligations.
  \item \texttt{z3\_binders: x: int, y: int} --- explicit Z3 sorts
    overriding the typed-Z3 encoder's inference.
  \item \texttt{cubical\_dim: 2} / \texttt{cohomology\_level: 1} ---
    request specific dimension / level for the cubical analysis.
  \item \texttt{expects\_sha256: "<hex>"} --- a commitment to the
    inspected source bytes.  The pipeline hashes the live
    \texttt{inspect.getsource} output; mismatch fires
    \texttt{hash\_drift\_count} and \texttt{certify\_or\_die}
    refuses certification.
  \item \texttt{expecting\_counterexample: true} --- assert that the
    counterexample search \emph{should} produce one (i.e.\ the
    property is supposed to be false).  Useful for testing the
    counterexample finder; if no counterexample is produced when one
    is expected, certification fails.
\end{itemize}

\section{Lemmas}
\label{sec:lemmas}

A \texttt{lemma} declaration introduces an intermediate fact:

\begin{lstlisting}[language=deppyfile]
lemma triangle_inequality_aux:
  statement: "(a + b)**2 <= 2 * (a**2 + b**2) when a >= 0 and b >= 0"
  lean_proof: |
    intros a b ha hb
    nlinarith [sq_nonneg (a - b)]
\end{lstlisting}

A lemma can be either:
\begin{itemize}[leftmargin=1.7em]
  \item \emph{admitted} (no \texttt{lean\_proof}; certificate emits
    \texttt{axiom Lemma\_X : Lemma\_X\_prop});
  \item \emph{proved by user-supplied \Lean{}}; the certificate
    emits \texttt{theorem Lemma\_X : Lemma\_X\_prop := by ...}.
\end{itemize}

Either way, the lemma is registered as a kernel axiom (file:
\texttt{deppy/lean/sidecar\_certificate.py}, lines around 183--188:
``\textbf{AUDIT 1.8}: register every parsed lemma as a kernel
axiom'').  Subsequent \texttt{verify} blocks may cite it via
\texttt{apply(triangle\_inequality\_aux)} in PSDL.

\section{The \texttt{code\_types:} block}
\label{sec:code-types}

A new top-level block lets the user provide typed signatures for
helper symbols PSDL references.

\begin{lstlisting}[language=deppyfile,caption={A \texttt{code\_types:} block.}]
code_types: |
  sqrt: Int → Int
  sum_zip_sub_sq: Point → Point → Int
\end{lstlisting}

These drive concrete \texttt{axiom}s in the Lean preamble:
\begin{lstlisting}[language=lean4]
axiom sqrt : Int → Int
axiom sum_zip_sub_sq : Point → Point → Int
\end{lstlisting}

The point is that \PSDL{} can mention \texttt{sqrt(...)} and
\texttt{sum\_zip\_sub\_sq(self, other)} in the user's natural
notation, and the certificate's preamble will declare these as
typed Lean symbols rather than opaque Int constants.  Without
\texttt{code\_types}, every helper symbol becomes
\texttt{opaque h : Int} and the surrounding tactics fail to
elaborate.

\section{Top-level \Lean{} integration}
\label{sec:lean-toplevel}

Two top-level keys inject text directly into the certificate:

\begin{itemize}[leftmargin=1.7em]
  \item \texttt{lean\_imports: |} --- one \texttt{import} per line,
    inserted at the top of the generated certificate.
  \item \texttt{lean\_preamble: |} --- raw \Lean{} text inserted
    after the imports and before the annotation-driven preamble.
\end{itemize}

These are the escape hatches.  When the user's proofs require
particular \Lean{} machinery (a Mathlib-style import, a custom
\texttt{macro}), the \texttt{.deppy} file can introduce it.  The
\texttt{certify\_or\_die} gate does not refuse certification on the
basis of which imports were used --- only on whether \Lean{} accepts
the resulting certificate.

\section{The legacy \texttt{proof} form}
\label{sec:legacy-proofs}

The original \texttt{.deppy} format used \texttt{proof}
declarations:

\begin{lstlisting}[language=deppyfile,caption={Legacy \texttt{proof} form
from the SymPy sidecar.}]
proof distance_nonneg_proof:
  target: sympy.geometry.point.Point
  by: foundation Real_sqrt_nonneg

proof distance_pythagoras_proof:
  target: sympy.geometry.point.Point
  by: rewrite Real_sqrt_sq_nonneg then cite dist_pythagoras

proof distance_zero_iff_equal_proof:
  target: sympy.geometry.point.Point
  by: chain Real_sqrt_zero_iff_zero, Real_sum_of_nonneg_zero_iff_all_zero, Real_sub_sq_zero_iff_eq

proof distance_is_metric_proof:
  target: sympy.geometry.point.Point
  by: composition dist_nonneg & dist_zero_iff_equal & dist_symmetric & dist_triangle
\end{lstlisting}

The \texttt{by:} clause supports five forms:
\begin{itemize}[leftmargin=1.7em]
  \item \texttt{by: foundation NAME} --- proof body is
    \verb|exact Foundation_NAME|;
  \item \texttt{by: axiom NAME} --- \verb|exact Sidecar_NAME|;
  \item \texttt{by: refl} --- \verb|rfl|;
  \item \texttt{by: chain F1, F2, F3} --- \verb|F1.trans F2 |.\verb|trans F3|;
  \item \texttt{by: rewrite F then cite A} --- \verb|rw [F]; exact A|;
  \item \texttt{by: composition F1 \& F2 \& F3} --- \verb|And.intro F1 (And.intro F2 F3)|.
\end{itemize}

The \texttt{verify} form supersedes \texttt{proof} for new code: it
exposes more knobs (binders, requires/ensures, PSDL, by\_lean) and
its trust accounting is more granular.  We retain the legacy form
for backwards compatibility.

\section{File layout reference}
\label{sec:layout-ref}

A complete sidecar file has roughly the following layout:

\begin{lstlisting}[language=deppyfile,caption={Canonical layout of a sidecar file.}]
# 1. Header
module: <package>
version: ">=<min>"

# 2. Top-level Lean integration (optional)
lean_imports: |
  import Mathlib.Tactic
  import Std

lean_preamble: |
  -- raw Lean text injected verbatim

# 3. Helper signatures (optional)
code_types: |
  helper1: T1 → T2
  helper2: T1 → T2 → T3

# 4. Foundations (the trust surface)
foundation F1: "stmt1"
foundation F2: "stmt2 when guard"
foundation F3:
  statement: "stmt3"
  citation: "ref"
  lean_proof: |
    <Lean tactics>

# 5. Specs (per class/function)
spec <dotted.path>:
  guarantee: "..."
  pure: true
  axiom A1: "stmt"
  axiom A2: "stmt"

# 6. Lemmas (intermediate facts)
lemma L1:
  statement: "..."
  lean_proof: |
    <Lean tactics>

# 7. Verifies (units of verification)
verify V1:
  function: <dotted.path>
  property: "..."
  binders: x: T, y: U
  via: foundation F1
  psdl_proof: |
    <PSDL tactics>

# 8. Legacy proofs (optional)
proof P1:
  target: <dotted.path>
  by: chain F1, F2
\end{lstlisting}

\section{The \texttt{SidecarModule} object}
\label{sec:sidecar-module}

The parsed file is a Python object with the following relevant
fields (from \texttt{deppy/proofs/sidecar.py}):

\begin{lstlisting}[caption={\texttt{SidecarModule} excerpt --- the
key fields and their types.}]
class SidecarModule:
    name: str
    target_module: str
    version: str
    _specs: dict[str, ExternalSpec]
    _axioms: dict[str, AxiomDecl]
    _foundations: dict[str, FoundationDecl]
    _proofs: list[ProofDecl]
    _verifies: list[VerifyDecl]
    _lemmas: list[LemmaDecl]
    _foundation_deps: dict[str, list[str]]
    _imports: list[str]
    _symbolic_constants: dict[str, str]
    _predicates: dict[str, str]
    _code_types: dict[str, str]      # CODE-TYPES — typed signatures
    _lean_imports: list[str]
    _lean_preamble: str
    _about_specs: list[tuple[Any, _SpecMetadata]]
\end{lstlisting}

The pipeline is library-agnostic: every transformation operates on
this dataclass.  The transformations are documented in
\Cref{ch:pipeline}.

\section{Worked example: \texttt{sympy\_geometry.deppy}}

We end this chapter by walking through the SymPy sidecar at the top
of the repository.  The file is 508 lines and breaks into these
sections:

\begin{enumerate}[leftmargin=1.7em]
  \item Lines 1--46: comments and header (\texttt{module: sympy},
    \texttt{code\_types: sqrt: Int → Int} and
    \texttt{sum\_zip\_sub\_sq: Point → Point → Int}).
  \item Lines 48--65: 14 \texttt{foundation} declarations
    (\texttt{Real\_sqrt\_nonneg}, \texttt{Real\_add\_comm},
    \dots).
  \item Lines 70--82: \texttt{spec sympy.geometry.point.Point} with
    8 \texttt{axiom}s about distance and collinearity.
  \item Lines 89--96: \texttt{spec sympy.geometry.line.Segment} with
    5 \texttt{axiom}s about midpoint and length.
  \item Lines 105--120: \texttt{spec sympy.geometry.polygon.Triangle}
    with 16 \texttt{axiom}s about centroid, area, Euler line, etc.
  \item Lines 127--134: \texttt{spec sympy.geometry.ellipse.Circle}
    with 5 \texttt{axiom}s.
  \item Lines 142--149: \texttt{spec sympy.geometry.polygon.Polygon}
    with 5 \texttt{axiom}s.
  \item Lines 161--393: 38 legacy \texttt{proof} declarations.
  \item Lines 404--508: 9 \texttt{verify} blocks, including the
    PSDL example \texttt{dist\_nonneg\_psdl}.
\end{enumerate}

The verify-block list:
\begin{itemize}[leftmargin=1.7em]
  \item \texttt{dist\_nonneg\_code} --- via foundation
    \texttt{Real\_sqrt\_nonneg}, with fibration descent.
  \item \texttt{dist\_pythagoras\_code} --- via foundation
    \texttt{Real\_sqrt\_sq\_nonneg}.
  \item \texttt{dist\_symmetric\_code} --- via
    \texttt{Real\_sub\_sq\_swap}.
  \item \texttt{segment\_contains\_code}.
  \item \texttt{triangle\_area\_nonneg\_code}.
  \item \texttt{triangle\_perimeter\_sum\_code}.
  \item \texttt{triangle\_centroid\_avg\_code}.
  \item \texttt{segment\_length\_code}.
  \item \texttt{polygon\_area\_nonneg\_code}.
  \item \texttt{circle\_radius\_eq\_code}.
  \item \texttt{dist\_nonneg\_psdl} --- the worked PSDL example.
\end{itemize}

The \texttt{dist\_nonneg\_psdl} block (lines 487--508) is the
canonical \PSDL{} demonstration: \texttt{psdl\_proof} contains a
real Python-syntax proof script.  We will return to this block
repeatedly throughout the rest of the book.

\section{Summary}
\label{sec:ch5-summary}

The \texttt{.deppy} sidecar format:
\begin{itemize}[leftmargin=1.7em]
  \item names a target Python module;
  \item lists \emph{foundations} as the trust surface;
  \item attaches \emph{specs} with \emph{axioms} per target;
  \item declares \emph{lemmas} for re-use;
  \item drives the verifier via \emph{verify} blocks with optional
    \PSDL{} or raw \Lean{} proof scripts;
  \item integrates with \Lean{} via \texttt{lean\_imports} and
    \texttt{lean\_preamble};
  \item supports a legacy \texttt{proof} form for back-compat.
\end{itemize}

The next chapter introduces \PSDL: the Python-syntax tactic
language that fills the \texttt{psdl\_proof:} field.
